\documentclass[11pt]{article}
\usepackage{amsmath,amssymb,amsthm,hyperref}
\newtheorem{theorem}{Theorem}
\newtheorem{lemma}[theorem]{Lemma}
\newtheorem{proposition}[theorem]{Proposition}
\newtheorem{remark}[theorem]{Remark}
\begin{document}
\title{Logarithmic tangent spaces and factorially many positive critical points}
\author{Research checkpoint, unchanged HeadComplexity model}
\date{24 November 2026}
\maketitle

\paragraph{Verification and scope.}
The identities and induction below are complete written proofs. They have not
been formalized in Lean. The companion exact-arithmetic checker verifies the
explicit base, the three-dimensional formula, and additional finite rational
instances. No numerical feasibility assertion is used in the induction.
This note proves neither of the accepted asymptotic targets. It rules out one
specific proposed estimate on positive-chamber critical points. It also records
why such critical-point multiplicity alone does not establish sign-pattern
richness.

\section{Exact tangent-space formulation}

Let $X=\{0,1\}^n$, and let $v(x)=(1,x_1,\ldots,x_n)$. For an affine form
$B_\beta(x)=\beta\cdot v(x)$ positive on $X$, define the vector
$\ell(\beta)=(\log B_\beta(x))_{x\in X}$. Work in any open parameter cone
where the denominator has strict original admissible slope signs. For $H$
denominators put
\[
 \Phi(\beta_1,\ldots,\beta_H)=\sum_{h=1}^H\ell(\beta_h).
\]
Its differential is exactly
\[
 D\Phi(\beta)[a](x)=\sum_{h=1}^H
       \frac{a_h\cdot v(x)}{B_{\beta_h}(x)}.
\]
The affine numerator directions $a_h$ are arbitrary. A constant bias is also
included, since replacing $a_1\cdot v$ by $a_1\cdot v+cB_{\beta_1}$ adds $c$.
Thus, for $H\geq1$, every original score with $H$ positive affine denominators
is a tangent vector of $\Phi$, and conversely. This statement does not assert
that the derivative parameters themselves are denominator parameters; only a
sufficiently small perturbation is needed to remain in the open cone.

For a vector $\lambda\in\mathbb R^X$, define the scalar master function
\[
 \Psi_\lambda(\beta)=\sum_{x\in X}\lambda_x\log B_\beta(x).
\]
Then $\lambda$ annihilates the complete numerator tangent space at
$\beta_1,\ldots,\beta_H$ if and only if every $\beta_h$ is a critical point of
the same master function. Indeed, the respective equations are
\[
 \sum_{x\in X}\frac{\lambda_xv_i(x)}{B_{\beta_h}(x)}=0
 \qquad(0\leq i\leq n,\ 1\leq h\leq H).
\]
Positive $\lambda_x$ would give concavity of the master function, but the
annihilator equations require mixed weights. Positive-weight uniqueness
results therefore do not apply to this reduction.

\section{A factorial positive critical-point family}

We normalize the constant denominator coefficient to one. For
$b=(b_1,\ldots,b_n)\in(0,\infty)^n$ and weights
$\lambda=(\lambda_1,\ldots,\lambda_n)$, put
\[
 P^{(n)}_\lambda(b)=
 \sum_{\varnothing\ne S\subseteq[n]}
       \lambda_{|S|}\log\left(1+\sum_{i\in S}b_i\right).
\]
This function is symmetric in the coordinates of $b$. Let $M_n(b)$ be its
gradient matrix with respect to the weights:
\[
 (M_n(b))_{ik}=
 \sum_{\substack{S\subseteq[n],\ |S|=k\\ i\in S}}
        \frac{1}{1+\sum_{j\in S}b_j},
 \qquad \nabla P^{(n)}_\lambda(b)=M_n(b)\lambda.
\]
Also define
\[
 q_n(\lambda)=\sum_{k=1}^n\binom{n-1}{k-1}\lambda_k.
\]

\begin{theorem}
For every integer $n\geq2$ there exist rational positive coordinates
\[
 b_1=b_2=1<b_3<\cdots<b_n
\]
and rational weights $\lambda_1=2,\lambda_2,\ldots,\lambda_n$ such that:
\begin{enumerate}
\item $M_n(b)\lambda=0$ and $\operatorname{rank}M_n(b)=n-1$;
\item the Hessian of $P^{(n)}_\lambda$ at $b$ is nonsingular;
\item $q_n(\lambda)\ne0$.
\end{enumerate}
Consequently, one such master function has at least $n!/2$ distinct
nondegenerate critical points in the strict positive orthant.
\end{theorem}

\begin{proof}
For $n=2$ take $b=(1,1)$ and $\lambda=(2,-3)$. The gradient vanishes,
the two rows of $M_2$ are $(1/2,1/3)$, and
\[
 \operatorname{Hess}P^{(2)}_\lambda(1,1)=
 \begin{pmatrix}-1/6&1/3\\1/3&-1/6\end{pmatrix},
 \quad \det\operatorname{Hess}P^{(2)}_\lambda=-1/12,
 \quad q_2(\lambda)=-1.
\]

Suppose the three properties hold in dimension $n$. Append a coordinate
$b_{n+1}=T$, and let the new weights be
$\mu(T)=(\mu_1(T),\ldots,\mu_{n+1}(T))$, with $\mu_1(T)=2$.
Keep one of the identical first two old gradient rows, all the other old
rows, and the new gradient row. Multiply the new row by $T$. These are $n$
linear equations for $\mu_2,\ldots,\mu_{n+1}$.

At $u=1/T=0$, the old $n-1$ independent rows are those of $M_n(b)$ and have
zero last entry. Their columns $2,\ldots,n$ form an invertible square
matrix: otherwise a vector in the kernel with first coordinate zero would
contradict the one-dimensional kernel spanned by $\lambda$, whose first
coordinate is $2$. The limiting last row is
\[
 \left(\binom n0,\binom n1,\ldots,\binom nn\right),
\]
whose last entry is one. Therefore the limiting square coefficient matrix
is invertible. The solution $\mu(T)$ exists for all sufficiently large $T$
and is rational analytic in $u$ near zero. It satisfies
\[
 (\mu_1(T),\ldots,\mu_n(T))=\lambda+O(T^{-1}),
 \qquad
 L(\mu(\infty))=0,
\]
where
\[
 L(\mu)=\sum_{k=1}^{n+1}\binom n{k-1}\mu_k.
\]
All $n+1$ gradient equations vanish, because the first two rows remain
identical. The gradient matrix has rank $n$: the selected $n$ rows are
independent for sufficiently large $T$, and the repeated row bounds its
rank by $n$.

Write $s=\sum_{i=1}^n b_i>0$ and
\[
 A(\mu)=\sum_{k=2}^{n+1}\binom{n-1}{k-2}\mu_k.
\]
Pascal's identity gives
\[
 L(\mu)=A(\mu)+q_n(\mu_1,\ldots,\mu_n),
 \qquad A(\mu(\infty))=-q_n(\lambda)=:a\ne0.
\]
For $S\subseteq[n]$, let $d_S=1+\sum_{i\in S}b_i$. The new-coordinate
gradient equation and two finite combinatorial sums are
\[
 0=\sum_{S\subseteq[n]}\frac{\mu_{|S|+1}(T)}{T+d_S},
 \qquad
 \sum_S\mu_{|S|+1}=L(\mu),
 \qquad
 \sum_S d_S\mu_{|S|+1}=L(\mu)+sA(\mu).
\]
Expanding the first identity at $T=\infty$ is legitimate because it is a
finite sum and $\mu(T)$ is analytic and bounded. It yields
\[
 0=\frac{L(\mu(T))}{T}
     -\frac{L(\mu(T))+sA(\mu(T))}{T^2}+O(T^{-3}),
 \qquad
 L(\mu(T))=\frac{sa}{T}+O(T^{-2}).
\]
In particular $q_{n+1}(\mu(T))=L(\mu(T))\ne0$ for all sufficiently large
$T$.

For the new Hessian entry, the same expansion gives
\[
 \begin{aligned}
 (\operatorname{Hess}P^{(n+1)}_{\mu(T)})_{n+1,n+1}
 &=-\sum_S\frac{\mu_{|S|+1}(T)}{(T+d_S)^2}\\
 &=-\frac{L(\mu(T))}{T^2}
    +\frac{2(L(\mu(T))+sA(\mu(T)))}{T^3}+O(T^{-4})\\
 &=\frac{sa}{T^3}+O(T^{-4}).
 \end{aligned}
\]
The old Hessian block converges to the nonsingular old Hessian; every
cross entry involving the new coordinate is $O(T^{-2})$. Hence the Schur
complement of the old block is
\[
 \frac{sa}{T^3}+O(T^{-4})\ne0
\]
for sufficiently large $T$. This proves nonsingularity of the full Hessian.
Choose an integer $T$ sufficiently large also to exceed all old coordinates.
The new coordinates and weights are rational and satisfy every inductive
property.

Finally, permutation symmetry carries the critical point to all its
coordinate permutations and carries its Hessian to a permutation-conjugate
matrix. Exactly two coordinates coincide, so the orbit has $n!/2$ points.
All remain in the strict positive orthant.
\end{proof}

\paragraph{Explicit three-dimensional check.}
The first induction step can be solved exactly for every $T>1$:
\[
 b=(1,1,T),\qquad
 \lambda=\left(2,-\frac{3(T+2)}{T+1},
                    \frac{4(T+3)}{T+1}\right),
 \qquad q_3(\lambda)=\frac{2}{T+1},
\]
\[
 \det\operatorname{Hess}P^{(3)}_\lambda(b)
 =-\frac{(T-1)^3(T+4)}{6(T+1)^4(T+2)^2(T+3)}.
\]
At $T=2$, rescaling the weights to make $\lambda_3=1$ gives
$(3/10,-3/5,1)$, with three distinct positive regular critical points.

\section{Full affine annihilators and the exact limitation}

The normalized critical points above also produce full affine annihilators.
Set $\Lambda_S=\lambda_{|S|}$ for nonempty $S$ and
\[
 \Lambda_\varnothing=-\sum_{\varnothing\ne S\subseteq[n]}\Lambda_S.
\]
At any normalized critical point $b$, contraction of its gradient with $b$
gives
\[
 0=\sum_{S\ne\varnothing}\Lambda_S
       \frac{\sum_{i\in S}b_i}{1+\sum_{i\in S}b_i}
  =\sum_{S\ne\varnothing}\Lambda_S
    -\sum_{S\ne\varnothing}\frac{\Lambda_S}{1+\sum_{i\in S}b_i}.
\]
Therefore $\sum_S\Lambda_S/B_b(S)=0$ as well as all $n$ slope equations.
The same $\Lambda$ annihilates every affine-numerator space at every point
of the permutation orbit. All denominators have positive constant term
and strictly positive slopes, so they are fully positive on the cube and
have the original monotone affine admissibility. Scaling the weights does
not change this fact.

The factorial count rules out an absolute bound $C^n$ for all mixed-weight
positive-chamber critical points of this normalized Boolean master family.
It does not rule out the accepted matching lower bound. In particular, for
one denominator $B>0$,
\[
 \operatorname{sgn}\frac{A(x)}{B(x)}=\operatorname{sgn}A(x).
\]
Thus the entire union of single-head tangent spaces has exactly the same
strict sign patterns as affine functions. The many critical branches do
not increase single-head sign capacity. Even selecting many of the above
denominators leaves their sum of numerator tangent spaces inside the same
hyperplane $\Lambda^\perp$. A useful counting argument must control images
or strict orthants, rather than count these parameter preimages.

\end{document}
